\section{The Hodge degree on the surfaces in the covolume estimate}
\label{sec:hodge-surface-degree}

The local rank estimates control the finite indices. To estimate the row covolume one also needs a Hodge intersection on an arithmetic surface, uniformly as the level and rank vary. For a principally polarized
abelian scheme we compare the second arithmetic Chern
character of the Hodge bundle with the degree of its
determinant on the generic curve. The uniform proportionality
is the surface input to the flag and component calculations
that follow. For the abelian scheme \(\pi:A\to S\) with zero section \(e\), put \(\omega=e^*\Omega^1_{A/S}\) and \(D=\det\omega\). Write \(\bar\omega\) for the Hodge bundle with its Hodge metric and \(d_D\) for the degree of \(D\) on the generic curve over \(\mathbb Q\). The coefficient \(r_*\) in Theorem~\ref{thm:hodge-surface-degree} is independent of the base and of the abelian scheme.

\begin{construction}
\[
\begin{array}{lll}
\text{base surface}&\text{degree or comparison}&\text{value}\\ \hline
\text{regular, flat over }\mathbb Z&
 \widehat{\deg}\,\widehat{\operatorname{ch}}_2(\bar\omega)&r_*d_D/2\\
\text{flat, computed on a cover }S'&
 \widehat{\deg}_{S'}\widehat{\operatorname{ch}}_2(\bar\omega_{S'})/[S':S]&r_*d_D/2\\
\text{over a field of characteristic }p&
 \deg\operatorname{ch}_2(\omega)&0
\end{array}
\]
\caption{The base and the degree used in the second Hodge intersection. The metric and arithmetic degree occur on the horizontal branch.}\label{fig:hodge-second-surfaces}
\end{construction}

\begin{theorem}[The second Hodge degree]\label{thm:hodge-surface-degree}
Let \(\pi:A\to S\) be a principally polarized abelian scheme of
relative dimension \(g\) over an integral projective arithmetic
surface, flat over the integers. Write \(\omega=e^*\Omega^1_{A/S}\)
and \(D=\det\omega\). Denote the geometric degree of \(D\) on
the generic curve by \(d_D\), with degree taken over \(\Q\).
Over the complex fibres use the Hodge metric from the Riemann
form of the polarization. There is a real constant \(r_*\), independent of the polarized
abelian scheme and of the base, such that
\begin{equation}\label{eq:hodge-surface-degree}
   \widehat{\deg}\,\widehat{\operatorname{ch}}_2(\bar\omega)
                  =\tfrac12 r_*d_D .
\end{equation}
The same constant must work as the level, base, polarization
and relative dimension vary.

The arithmetic degree here is
the one on a regular model or, for a general surface in the
construction, the degree on a regular generically finite cover
divided by the degree of that cover. The flag description uses
half the sum of the arithmetic squares of the graded lines with
the Bott--Chern correction for their metrics. This definition
agrees on dominating regular covers by the projection formula
and functoriality of the secondary forms. Over an integral projective surface in characteristic \(p>0\),
the Hodge bundle of a principally polarized abelian scheme
satisfies \(\deg\operatorname{ch}_2(\omega)=0\). \label{cor:hodge-surface-degree-1}
\end{theorem}

\begin{proof}
\proofstep{Arithmetic Riemann--Roch on a regular theta model.}

First take \(S\) regular and a symmetric rigidified ample theta
line \(L\) inducing the principal polarization. Give \(L\) the
cubical metric and use on \(\omega\) the flat-family convention
for arithmetic Riemann--Roch. In a period chart this convention
uses the dual Riemann metric with parameter \(1/2\).
The cubical normalization makes \([2]^*\bar L\) isometric to
\(\bar L^4\). One obtains the metric by correcting the doubling
error of a chosen metric by the convergent series in iterated
doubling with coefficients \(4^{-j}\). In period coordinates it
is the translated Gaussian theta metric, hence smooth. On a geometric fibre \(A_s\), ample vanishing and the degree of
the principal polarization give
\[
 H^i(A_s,L_s)=0\ (i>0),\qquad
 h^0(A_s,L_s)=\chi(A_s,L_s)=\deg(c_1(L_s)^g)/g!=1 .
\]
Cohomology and base change therefore gives \(R^i\pi_*L=0\) for \(i>0\)
and a line \(M=\pi_*L\) carrying the fibre \(L^2\) metric. Use the
induced metric on \(D\).

Write \(a\) for the arithmetic second Chern character degree in
this metric. There are two terms from analytic torsion on the
complex base curve. Write \(t\) for its degree-zero term and
\(I\) for the integrated degree-two term, in the conventions
below. Pairings of metrized lines on \(S\) denote arithmetic
intersection degrees. For any smooth metrized line \(N\) on
\(S\), write \(d_N\) for its geometric degree. Arithmetic
Riemann--Roch gives
\begin{align}
 (M,N) - \tfrac t2 d_N &= -\tfrac12(D,N),\label{eq:hodge-arr-linear}\\
 \tfrac12(M,M)-\tfrac I2
     &=\tfrac18(D,D)-\tfrac{a}{12}+\tfrac{R_1}{2}d_D ,
     \label{eq:hodge-arr-square}
\end{align}
where \(R_1=2\zeta'(-1)+\zeta(-1)\) is the linear coefficient
of the arithmetic \(R\)-genus. Here the data to which Riemann--Roch is applied
are a smooth projective abelian scheme over a regular arithmetic
surface, its principal theta bundle, the polarization's
translation-invariant relative Kähler metric and the \(L^2\)
metric on the direct image. The higher-degree theorem includes
the Bismut--Köhler analytic torsion form in this direct image.
\cite{arithmetic-rr-surface}\label{cor:hodge-surface-degree-2}

For the coefficients in these equations, expand the Todd
class for \(\pi^*\omega^\vee\) in degrees one and two. For Chern roots \(v_i\) of \(\omega\), the dual has roots \(-v_i\).
Multiplying the individual Todd factors to degree two gives
\[
 \begin{aligned}
 \operatorname{Td}_1&=-\tfrac12\sum_i v_i,\\
 \operatorname{Td}_2&=\tfrac1{12}\sum_i v_i^2+\tfrac14\sum_{i<j}v_iv_j
 =\tfrac18\Big(\sum_i v_i\Big)^2-\tfrac1{24}\sum_i v_i^2 .
 \end{aligned}
\]
For the arithmetic classes the second sum includes the flag
metric correction. The correction from \(1-R\) uses the
linear term on the dual, giving \(+R_1d_D/2\) in degree.
The degree-\(g\) power of \(c_1(L)\), divided by \(g!\),
has fibre degree one.

For any base class \(B\) of complementary degree, the projection
formula and the cubical metric give
\[
 4^g\,\widehat{\deg}\bigl(c_1(\bar L)^k\pi^*B\bigr)
  =\widehat{\deg}\bigl([2]^*(c_1(\bar L)^k\pi^*B)\bigr)
  =4^k\,\widehat{\deg}\bigl(c_1(\bar L)^k\pi^*B\bigr).
\]
The first factor is the degree of \([2]\) on the abelian scheme.
Hence the terms with \(k=g+1,g+2\) vanish. Pairing the degree-one formula with
\(N\) and taking the degree of the degree-two formula gives
\eqref{eq:hodge-arr-linear} and \eqref{eq:hodge-arr-square}.

\proofstep{The two analytic torsion terms.}
Here \(t=t_g\) is constant and \(I=8\pi l_g d_D\) for a real
\(l_g\) depending only on \(g\). We explain the invariance
that removes the arithmetic model from these two coefficients.
Over a simply connected period chart write the underlying
real torus as a product using flat period coordinates \(t\).
The Chern form \(\Omega_L\) on the total space satisfies
\([2]^*\Omega_L=4\Omega_L\) by the cubical isometry.

In these
coordinates doubling is \((b,t)\mapsto(b,2t)\), with \(b\) the
base variable. A coefficient with \(j\) vertical differentials
satisfies
\[
 2^j c(b,2t)=4c(b,t).
\]
For fixed \(b\), put \(M_b=\sup_t|c(b,t)|\), finite by compactness.
Iteration for each \(j<2\) gives
\[
 |c(b,t)|=2^{-h(2-j)}|c(b,2^h t)|
 \le 2^{-h(2-j)}M_b\longrightarrow0 .
\]
Thus only the purely vertical term survives. The fibre polarization prescribes the component
with two vertical differentials. Hence the full form is the
constant polarization form in the flat real coordinates.

Two theta lines with the same total curvature differ by a flat
unitary line. On the simply connected base chart such a line is
determined by its character \(\xi:\Lambda\to\mathrm U(1)\) on the
period lattice. If \(E_L\) is the alternating polarization form,
the holonomy ratio for translation by a flat vector \(u\) is
\[
 \operatorname{hol}_{t_u^*\bar L}(\lambda)
   \operatorname{hol}_{\bar L}(\lambda)^{-1}
      =\exp(2\pi i E_L(u,\lambda)).
\]
This follows by integrating the curvature over the cylinder swept
by \(\lambda\). Nondegeneracy of \(E_L\) lets us prescribe phases
on a basis of \(\Lambda\). Choose \(u\) giving \(\xi\). The
translated holomorphic datum has the same horizontal splitting
and relative metric. The remaining flat quotient has trivial
holonomy and its parallel unit section is the desired isometry.

For \(\Lambda'\subset\Lambda\) of finite index, push the forms
down along the cover with deck group \(A=\Lambda/\Lambda'\).
The character decomposition gives the spaces for
\(\bar L\otimes F_\chi\), \(\chi\in\widehat A\), where \(F_\chi\)
is the unitary period twist. The horizontal differential and
vertical Laplacian preserve each character space. If
\(\mathcal T(\Lambda,\bar L)\) is the torsion form, additivity
of the superconnection trace and translation invariance give
\[
 \mathcal T(\Lambda',\bar L|_{\Lambda'})
   =\sum_{\chi\in\widehat A}\mathcal T(\Lambda,\bar L\otimes F_\chi)
   =[\Lambda:\Lambda']\,\mathcal T(\Lambda,\bar L).
\]
Here \(\bar L|_{\Lambda'}\) denotes pullback to the covering torus.
For a rational symplectic change of marking \(g\), the lattices
\(\Lambda\) and \(g\Lambda\) have the same covolume. Writing \(H=\Lambda\cap g\Lambda\), we have
\[
 [\Lambda:H]=\frac{\operatorname{covol}H}{\operatorname{covol}\Lambda}
 =\frac{\operatorname{covol}H}{\operatorname{covol}(g\Lambda)}
 =[g\Lambda:H].
\] Apply the last identity on the common
theta chart and cancel the index. This gives rational symplectic
invariance. Smooth dependence on the period and its tangent and
density give real symplectic invariance.
\label{cor:hodge-surface-degree-3}


This group acts transitively on Siegel space, so the scalar
part is constant. At \(iI_g\) its stabilizer \(U(g)\) acts on
the tangent by the symmetric square of the standard representation.
It is irreducible for \(g>0\). Thus the degree-two part is a
real multiple of the invariant Hermitian form. In the source
normalization its evaluation on the complex tangent \(v\) is
\(i l_g\,\mathrm{sdens}(\tau,v)\), where the density at
\(iI_g\) is \(\operatorname{tr}(v v^*)\).

For \(Y=\operatorname{Im}\tau\) the density at \(\tau\) is
\(\operatorname{tr}(Y^{-1}vY^{-1}v^*)\). Indeed differentiation
of \(-\log\det Y\) gives one quarter of this Hermitian form
as its complex Hessian:
\[
 -\partial_v\bar\partial_v\log\det Y
      =\tfrac14\operatorname{tr}(Y^{-1}vY^{-1}v^*).
\]
In the arithmetic integral the real area densities of the torsion
and the first Chern form of \(D\) are \(2l_g\,\mathrm{sdens}\)
and \(\mathrm{sdens}/(4\pi)\), respectively. The former density
is \(8\pi l_g\) times the latter, so \(I=8\pi l_gd_D\).
In dimension zero both degrees of the Hodge bundle vanish.
\label{cor:hodge-surface-degree-4}

\proofstep{Elimination of the direct-image intersections.}
Test \eqref{eq:hodge-arr-linear} first with the trivial line
of constant weight two. Its geometric degree is zero and its
arithmetic pairing with any line is that line's geometric
degree. Hence \(d_M=-d_D/2\). Testing with \(D\) and \(M\)
and using symmetry gives
\[
       (M,M)=\tfrac14(D,D)-\tfrac t2 d_D .
\]
Substitute in \eqref{eq:hodge-arr-square}. The self-intersection
of \(D\) cancels and
\[
       a=6I+(6R_1+3t)d_D
        =(48\pi l_g+6R_1+3t_g)d_D .
\]
To return to the original Hodge metric, put \(s=\log(2\pi)\).
Write \(a_{\rm H}\) for the arithmetic second Chern character degree
in the original Hodge metric. The metric used to compute \(a\)
multiplies the original squared metric by \(e^s\). The
Bott--Chern forms of the flag are invariant under this common
scaling. If its graded lines have generic degrees \(d_i\), the
change in each half-square is \(-s d_i/2\). Consequently
\begin{equation}\label{eq:hodge-change-of-metric}
 a-a_{\rm H}=-\frac{s}{2}\sum_i d_i=-\frac{s}{2}d_D,
 \qquad a_{\rm H}=\frac{r_g}{2}d_D,\qquad
 r_g=2(48\pi l_g+6R_1+3t_g)+\log(2\pi).
\end{equation}
\label{cor:hodge-surface-degree-5}

\proofstep{Descent through a regular alteration.}
For a regular generically finite cover \(\rho:S'\to S\) of
degree \(n>0\), invariant differentials and their Hodge metrics
pull back isometrically. Write primes for the invariants of the pullback.
The projection formula on the generic fibre and its arithmetic
counterpart on a regular flag model give
\[
 d_{D'}=n d_D,\qquad a'_{\rm H}=n a_{\rm H},\qquad
 a_{\rm H}-\tfrac12 r_gd_D
       =\frac{1}{n}\bigl(a'_{\rm H}-\tfrac12 r_gd_{D'}\bigr).
\]
One can
therefore calculate on a regular flag model. A symmetric
theta representative exists on the generic fibre after finite
extension.

On a regular cover over that extension take its
divisor closure on the abelian scheme and rigidify at zero.
The possible vertical changes in the extension of the generic
line and its symmetry are pulled back from the base, since
the smooth abelian fibres are geometrically integral. The
rigidification removes these changes. The homomorphism to the
dual abelian scheme induced by the line agrees with the given
polarization on the dense generic fibre and hence everywhere.
Thus the same two degree formulas descend the regular-theta
calculation to the original polarized family.
\label{cor:hodge-surface-degree-6}

\proofstep{A constant for every relative dimension.}
To obtain a constant uniform in \(g\), over
\(\Q(\sqrt{29})\) set \(e=(5+\sqrt{29})/2\). Then \(e(e-5)=1\).
The integral elliptic curve
\[
 y^2+xy+e^2y=x^3
\]
has
discriminant \(-e^{10}\), a unit, and has the origin theta
divisor. After passing to a regular cover factoring through the
integers of this field, its constant pullback has \(d_D=0\)
and arithmetic second Chern character degree zero. Its
metrized Hodge line is pulled back from the arithmetic curve
of integers, on which an arithmetic square vanishes.
On the regular cover let \(A_0\) denote the pulled-back
constant elliptic scheme. The product polarization gives an
orthogonal direct sum
\[
 \omega_{A\times A_0}\simeq\omega_A\oplus\omega_{A_0}.
\]

Taking determinants and second arithmetic Chern characters gives
\[
 d_{\det\omega_{A\times A_0}}=d_{\det\omega_A},\qquad
 \widehat{\deg}\,\widehat{\operatorname{ch}}_2(\omega_{A\times A_0})
       =\widehat{\deg}\,\widehat{\operatorname{ch}}_2(\omega_A).
\]
The terms from \(\omega_{A_0}\) vanish as just shown.
Each product raises the relative dimension while preserving both degrees. This allows the constants from the fixed-dimension formula to be compared on the same nonzero degree.

If some example on a regular theta model has \(d_D\ne0\),
choose one of least relative dimension \(g_0\). For any
\(g\ge g_0\) raise this example to dimension \(g\) by the
preceding construction. Use the original Hodge metric throughout this comparison.
The product on the cover has degrees \(n a_{\rm H}\) and \(n d_D\).
Applying \eqref{eq:hodge-change-of-metric} in dimensions \(g_0\) and \(g\)
to the same pair of numbers gives
\begin{equation}\label{eq:hodge-dimension-cancellation}
 n r_{g_0}d_D=2n a_{\rm H}=r_g(n d_D),\qquad n d_D\ne0 .
\end{equation}
Cancelling \(n d_D\) in \eqref{eq:hodge-dimension-cancellation}
gives \(r_g=r_{g_0}\). Thus one can use \(r_*=r_{g_0}\) for all these dimensions. In smaller dimensions
\(d_D=0\) for all examples, and their fixed-dimension formula
already gives degree zero. If every dimension has this
property, take \(r_*=0\). This proves the required uniformity,
including for covers whose degrees grow with the level.
\label{cor:hodge-surface-degree-7}

\proofstep{The finite-characteristic calculation.}
Over an integral projective surface in characteristic
\(p\), write \(F\) for absolute Frobenius. The Hodge and
conjugate filtrations of \(H^1_{\rm dR}(A/S)\) have graded
bundles \((\omega,\omega^\vee)\) and
\((F^*\omega^\vee,F^*\omega)\) by principal polarization and
the Cartier isomorphism. More explicitly, for \(H=H^1_{\rm dR}(A/S)\) the two exact sequences are
\[
 0\to\omega\to H\to R^1\pi_*\mathcal O_A\to0,\qquad
 0\to F^*R^1\pi_*\mathcal O_A\to H\to F^*\omega\to0.
\]
The Hodge spectral sequence degenerates at \(E_1\) for an abelian scheme and its Hodge cohomology sheaves are locally free. Relative Cartier identifies the conjugate graded pieces with their Frobenius pullbacks, compatibly with base change. These assertions apply over an arbitrary base of characteristic \(p\) by Moonen and Wedhorn \cite[\S\S6.2--6.5, Example~6.4(1)]{moonen-wedhorn-fzip}. The principal polarization identifies \(R^1\pi_*\mathcal O_A=\operatorname{Lie}(A^\vee/S)\) with \(\omega^\vee\). Additivity in the two sequences therefore computes the same degree-two Chern character. 

On a flag
cover a line root \(x\) becomes \(-x\) under duality and \(px\)
under Frobenius pullback. Thus
\[
 \begin{aligned}
 \operatorname{ch}_2(H^1_{\mathrm{dR}})
   &=\operatorname{ch}_2(\omega)+\operatorname{ch}_2(\omega^\vee)
    =2\operatorname{ch}_2(\omega),\\
 \operatorname{ch}_2(H^1_{\mathrm{dR}})
   &=\operatorname{ch}_2(F^*\omega^\vee)+\operatorname{ch}_2(F^*\omega)
    =2p^2\operatorname{ch}_2(\omega).
 \end{aligned}
\]
Taking degrees gives zero since \(p^2-1\ne0\), including for
\(p=2\).
\end{proof}

In the de Rham comparison, the vertical terms contribute zero by this calculation and the horizontal terms use the common constant \(r_*\). The flag calculation in
Section~\ref{sec:covolume-flag-test} combines these degree
identities with its volume, semipositivity and \(p\)-adic
comparison inputs. Section~\ref{sec:flag-component-transfer}
applies the resulting estimate on the selected levels.
\label{cor:hodge-surface-degree-8}
